% Generated draft: requires mathematical and source review.
\documentclass[11pt]{article}
\usepackage[margin=25mm]{geometry}
\usepackage{fontspec}
\setmainfont{Latin Modern Roman}
\usepackage{amsmath,amssymb,mathtools,unicode-math}
\setmathfont{Latin Modern Math}
\usepackage{enumitem,needspace,xurl,hyperref}
\setlength{\parindent}{0pt}
\setlength{\parskip}{5pt}
\setlength{\emergencystretch}{3em}
\newcommand{\N}{\mathbb{N}}
\newcommand{\Z}{\mathbb{Z}}
\newcommand{\Q}{\mathbb{Q}}
\newcommand{\R}{\mathbb{R}}
\newcommand{\C}{\mathbb{C}}
\newcommand{\F}{\mathbb{F}}
\newcommand{\A}{\mathbb{A}}
\newcommand{\PP}{\mathbb{P}}
\newcommand{\E}{\mathbb{E}}
\newcommand{\eps}{\varepsilon}
\newcommand{\dd}{\,\mathrm{d}}
\begin{document}
\section*{UnsolvedMath ID: 701107}
\textbf{Source problem code:} catalogue entry
\subsection{Definitions and mathematical statement}
Let \(p\) be a prime, let \(X\) be an alphabet with \(|X|=p\), and let \(T_p=X^*\) be the rooted tree whose root is the empty word and whose edges join \(v\) to \(vx\) for \(x\in X\). Thus every vertex has exactly \(p\) children. For a root-preserving automorphism \(g\) and a vertex \(v\), its section \(g|_v\) is determined by
\[
 g(vw)=g(v)\,g|_v(w)\qquad(w\in X^*).
\]
An automorphism is finite-state if the set of its sections is finite; equivalently, it is represented by an invertible finite synchronous automaton. Under the literal tree-action wording of the supplied record, define
\[
\mathcal F_p=\{G: G\text{ is finitely generated and admits a faithful action on }T_p\text{ generated by finite-state automorphisms}\}.
\]
A group \(G\) is residually-\(p\) if, for every \(g\in G\setminus\{1\}\), there are a finite \(p\)-group \(P\) and a homomorphism \(\varphi:G\to P\) with \(\varphi(g)\ne1\).

\textbf{Source qualification.} The record asserts that every member of \(\mathcal F_p\) is residually-\(p\) and has word problem decidable in at most exponential time. It then states that residually-\(p\) groups with decidable word problem outside \(\mathcal F_p\) therefore exist, although no concrete example is known. These are supplied contextual assertions, including a historical knowledge claim, and are not independently verified here. In particular, for general primes the residual-\(p\) assertion does not follow from the literal definition above: the record does not specify any restriction on local permutations that would justify it. Whether such a restriction is intended is unresolved.

\textbf{Part (a): proposed examples.} The source refers to groups \(G_\omega\) in R. I. Grigorchuk, \emph{Math. USSR--Sb.} \textbf{54} (1986), 185--205. It asks whether some, or even all, of these groups fail to belong to \(\mathcal F_p\) when \(\omega\) is computable but not periodic. If \(\Omega\) denotes the admissible sequence domain of that cited construction, the two proposed assertions are
\[
\exists\omega\in\Omega\;\bigl(\omega\text{ is computable and not periodic}\ \land\ G_\omega\notin\mathcal F_p\bigr)
\]
and, more strongly,
\[
\forall\omega\in\Omega\;\bigl(\omega\text{ is computable and not periodic}\ \Longrightarrow\ G_\omega\notin\mathcal F_p\bigr).
\]
These displays express the source's existential and universal alternatives, subject to unresolved specifications: the record does not reproduce the construction of \(G_\omega\), the domain \(\Omega\), or its relationship to the prime \(p\). It also does not say whether ``not periodic'' excludes eventually periodic sequences. Thus part (a) is a source-dependent formulation requiring these details to be checked, rather than a complete intrinsic definition of the proposed groups.

\textbf{Part (b): a restricted wreath product.} For groups \(A,B\), use the restricted regular wreath product
\[
A\wr B=\left(\bigoplus_{b\in B}A\right)\rtimes B,
\]
where the direct sum consists of finitely supported functions \(B\to A\), with pointwise multiplication, and \(B\) acts by left translation: \((h\cdot f)(b)=f(h^{-1}b)\). Let \(\Z\) denote the infinite cyclic group. The question is
\[
\Z\wr(\Z\wr\Z)\ \stackrel{?}{\in}\ \mathcal F_2.
\]
Both wreath products are restricted. The supplied expression has malformed LaTeX braces; the displayed right-associated product is the reading indicated by its visible parentheses, and this transcription repair requires review.

\textbf{Supplied status.} The registry labels the overall problem partially solved. Its discussion reports an affirmative solution to part (b), attributed to A. C. Dantas, J. R. Oliveira, and T. M. G. Santos, preprint, 2024. This is a supplied literature report, not an independently verified theorem assessment. The record supplies no resolution of part (a).
\subsection{Short English statement}
For a prime \(p\), the problem asks whether some, or all, computable nonperiodic members of a cited Grigorchuk family fail to admit faithful finite-automaton actions on the rooted tree with \(p\) children per vertex. It also asks whether the restricted wreath product \(\Z\wr(\Z\wr\Z)\) admits such an action on the binary rooted tree; the supplied record reports an affirmative solution to this second part. The family parameters, meaning of nonperiodicity, and general-prime local-action convention require clarification.
\subsection{Sources}
\begin{itemize}
\item Supplied local UnsolvedMath registry record: permanent ID \(701107\), exact source problem code \texttt{catalogue entry}, title \emph{Kourovka Notebook Problem 15.19}. Proposed by L. Bartholdi and S. Sidki, 2002. The supplied registry status is \texttt{partially\_solved}.
\item \emph{The Kourovka Notebook (Unsolved Problems in Group Theory)}, problem 15.19, 15th issue (2002); supplied reference edition: 21st issue (2026), \texttt{arXiv:1401.0300v48}. The registry describes the collection as edited by E. I. Khukhro and V. D. Mazurov.
\item R. I. Grigorchuk, \emph{Math. USSR--Sb.} \textbf{54} (1986), 185--205. Supplied reference for the groups in part (a); no article title or construction is included in the record.
\item A. M. Brunner and S. Sidki, \emph{J. Algebra} \textbf{257} (2002), 51--64. Supplied background reference; no article title is included.
\item R. I. Grigorchuk, V. V. Nekrashevich, and V. I. Sushchanskiĭ, in \emph{Dynamical systems, automata, and infinite groups}, \emph{Proc. Steklov Inst. Math.} \textbf{231} (2000), 128--203. Supplied background reference.
\item A. C. Dantas, J. R. Oliveira, and T. M. G. Santos, preprint, 2024, \url{https://arxiv.org/pdf/2405.16678}. The supplied discussion attributes an affirmative solution of part (b) to this work; no title is supplied.
\item The supplied status review is dated \texttt{2026-10-16}, with evidence designated \texttt{maintained\_source}. It reports that part (b) is marked solved in the October 2026 21st issue, \texttt{arXiv:1401.0300v48}. This is supplied dated curation, not an independently verified assessment; no cited source has been visited for this draft.
\end{itemize}
\end{document}
